\begin{frame}{고정점과 그 ``속도''}
  위 계산을 10단계까지 이어가면 ($w_{xh}=0.5$, $w_{hh}=0.8$):
  \vskip0.4em
  \small
  \begin{tabular}{@{}l|cccccccc@{}}
    \toprule
    $t$ & 1 & 2 & 3 & 4 & 5 & 6 & 8 & 10 \\
    \midrule
    $z_t = 0.5 + 0.8\,h_{t-1}$ & 0.500 & 0.870 & 1.061 & 1.129 & 1.149 & 1.154 & 1.156 & 1.156 \\
    $h_t = \tanh(z_t)$ & 0.462 & 0.701 & 0.786 & 0.811 & 0.817 & 0.819 & 0.820 & 0.820 \\
    \bottomrule
  \end{tabular}
  \vskip0.5em
  \begin{itemize}
    \item $h_t$는 \kb{고정점(fixed point)}
      $h^\ast = \tanh(0.5 + 0.8\, h^\ast)$으로 수렴 --- 해는
      $h^\ast \approx 0.8196$, $t=10$에서 차이 $< 10^{-3}$
    \item ``같은 입력이 계속 오면 은닉 상태가 결국 한 곳에
      정지한다'' --- RNN의 \textbf{동적(dynamics)}을 이해하는
      출발점
    \item 수렴 \textbf{속도}는 고정점 근처의 ``\textbf{증폭 인자}''로
      결정: $\frac{d h_t}{d h_{t-1}} = 0.8 \cdot tanh'(z_t)$
    \item $h^\ast \approx 0.82$에서 $\tanh'(z^\ast) = 1 - (h^\ast)^2
      \approx 0.325$ $\Rightarrow$ 증폭 인자 $0.8 \times 0.325
      \approx 0.26$ --- 각 단계마다 차이가 약 4분의 1로 줄어듦
      (표에서 $h_5$ 이후가 거의 움직이지 않는 이유)
  \end{itemize}
\end{frame}
